\section{The Page That Got Slower}
\label{sec:ch18-the-page-that-got-slower}

\index{blame routing!the page nobody can attribute|(}%
Chapter~\ref{ch:ownership} ended on a graph that is slow, a request that
crossed several services, and a composed schema that says which service
answers each field and nothing about which one took the time. Here is that
graph. It is the four services this book has built, with one change I made
and did not announce: the Speakers service's reference resolver is the naive
one chapter~\ref{ch:router-n-plus-1} opened on, fetching one row per
representation instead of going through a DataLoader. I know which service it
is because I changed it. The person holding the pager does not, and everything
below is what that person has to work with.

\index{HTTP 200!a slower page looks identical}%
The complaint is about the session page, asked for with the two rating fields
chapter~\ref{ch:second-third-subgraph} contributed from the Ratings service, so
that one request reaches three of the four services:

\begin{minted}[fontsize=\footnotesize]{text}
POST /graphql HTTP/1.1
Host: localhost:3002
Content-Type: application/json

{"query":"{ sessions(first: 10) { nodes { title speaker { name } averageScore ratingCount } } }"}
\end{minted}

\index{router!response headers name no service}%
The body that comes back is byte for byte the one the graph returns without my
change: four sessions, every speaker, every rating there is, no \code{errors}
key. These
are the headers it arrives with, the date aside, and the thing to look for in
them is anything at all that names a service or identifies the request:

\begin{minted}{text}
HTTP/1.1 200 OK
Content-Type: application/json; charset=utf-8
Vary: Accept-Encoding
Content-Length: 463
\end{minted}

\index{StatementLog@\code{StatementLog}!a number and nothing else}%
Nothing does. So the next place to look is the logs, and there are five of
them. Four are the services' consoles, each printing
chapter~\ref{ch:data-without-n-plus-1}'s \code{StatementLog} markers. Here is
what the Speakers console printed while that request was being answered:

\begin{minted}{sql}
-- statement 1
SELECT "s"."Id", "s"."Bio", "s"."Email", "s"."Name"
FROM "Speakers" AS "s"
WHERE "s"."Id" = @key
LIMIT 1
-- statement 2
SELECT "s"."Id", "s"."Bio", "s"."Email", "s"."Name"
FROM "Speakers" AS "s"
WHERE "s"."Id" = @key
LIMIT 1
-- statement 3
SELECT "s"."Id", "s"."Bio", "s"."Email", "s"."Name"
FROM "Speakers" AS "s"
WHERE "s"."Id" = @key
LIMIT 1
\end{minted}

\index{N+1@N+1!invisible in a shared log}%
Three statements, one row each. Read alone, on a quiet laptop, that is the
answer, and chapter~\ref{ch:router-n-plus-1} spent a section explaining it.
On a service answering real traffic it answers nothing. The marker is a
process-wide counter and the text is a query. Neither line says which request
issued it, so three statements in a row could be one request paying three
times or three requests paying once each, and the Speakers team cannot tell
those apart from their own console any better than you can. The Sessions and
Ratings consoles printed one statement each, as a marker and a query in the
same form.

\index{router!the request log line}%
The fifth log is the router's. It writes one line per request, and this is
the line for the page, with its timestamp, its timing, the router's own
request number and the fields that identify this machine and this client
taken out:

\begin{minted}[fontsize=\footnotesize,breakanywhere]{text}
INFO requestlogger/requestlogger.go:206 /graphql {"log_type": "request", "method": "POST", "path": "/graphql", "config_version": "00000000-0000-0000-0000-000000000000", "trace_id": "1cd345017494455980971dceaf195769", "status": 200}
\end{minted}

\index{trace id!the only line that carries one}%
A \emph{trace id} is an identifier given to one request so that everything
done on its behalf, in any process, can be recognized as belonging to it, and
\code{trace\_id} on that line is one. It is the only identifier for this
request that appears in any of the five logs and was also sent anywhere
beyond the router, and the router is the only component that printed it. The Sessions, Speakers and Ratings consoles hold no line
containing those thirty-two characters, and the client that complained was
never shown them. So what the on-call engineer can prove is that the router
took a request and answered 200. Which of three services spent the effort is a
question with no key to join on.
\index{blame routing!the page nobody can attribute|)}%
